% Options for packages loaded elsewhere
\PassOptionsToPackage{unicode}{hyperref}
\PassOptionsToPackage{hyphens}{url}
\documentclass[
  ignorenonframetext,
]{beamer}
\newif\ifbibliography
\usepackage{pgfpages}
\setbeamertemplate{caption}[numbered]
\setbeamertemplate{caption label separator}{: }
\setbeamercolor{caption name}{fg=normal text.fg}
\beamertemplatenavigationsymbolsempty
% remove section numbering
\setbeamertemplate{part page}{
  \centering
  \begin{beamercolorbox}[sep=16pt,center]{part title}
    \usebeamerfont{part title}\insertpart\par
  \end{beamercolorbox}
}
\setbeamertemplate{section page}{
  \centering
  \begin{beamercolorbox}[sep=12pt,center]{section title}
    \usebeamerfont{section title}\insertsection\par
  \end{beamercolorbox}
}
\setbeamertemplate{subsection page}{
  \centering
  \begin{beamercolorbox}[sep=8pt,center]{subsection title}
    \usebeamerfont{subsection title}\insertsubsection\par
  \end{beamercolorbox}
}
% Prevent slide breaks in the middle of a paragraph
\widowpenalties 1 10000
\raggedbottom
\AtBeginPart{
  \frame{\partpage}
}
\AtBeginSection{
  \ifbibliography
  \else
    \frame{\sectionpage}
  \fi
}
\AtBeginSubsection{
  \frame{\subsectionpage}
}
\usepackage{iftex}
\ifPDFTeX
  \usepackage[T1]{fontenc}
  \usepackage[utf8]{inputenc}
  \usepackage{textcomp} % provide euro and other symbols
\else % if luatex or xetex
  \usepackage{unicode-math} % this also loads fontspec
  \defaultfontfeatures{Scale=MatchLowercase}
  \defaultfontfeatures[\rmfamily]{Ligatures=TeX,Scale=1}
\fi
\usepackage{lmodern}
\ifPDFTeX\else
  % xetex/luatex font selection
\fi
% Use upquote if available, for straight quotes in verbatim environments
\IfFileExists{upquote.sty}{\usepackage{upquote}}{}
\IfFileExists{microtype.sty}{% use microtype if available
  \usepackage[]{microtype}
  \UseMicrotypeSet[protrusion]{basicmath} % disable protrusion for tt fonts
}{}
\makeatletter
\@ifundefined{KOMAClassName}{% if non-KOMA class
  \IfFileExists{parskip.sty}{%
    \usepackage{parskip}
  }{% else
    \setlength{\parindent}{0pt}
    \setlength{\parskip}{6pt plus 2pt minus 1pt}}
}{% if KOMA class
  \KOMAoptions{parskip=half}}
\makeatother
\setlength{\emergencystretch}{3em} % prevent overfull lines
\providecommand{\tightlist}{%
  \setlength{\itemsep}{0pt}\setlength{\parskip}{0pt}}
\usepackage{bookmark}
\IfFileExists{xurl.sty}{\usepackage{xurl}}{} % add URL line breaks if available
\urlstyle{same}
\hypersetup{
  hidelinks,
  pdfcreator={LaTeX via pandoc}}

\author{\texorpdfstring{}{}}
\date{}

\begin{document}

\begin{frame}[fragile]{Infrastructure and connectivity}
\protect\phantomsection\label{infrastructure-and-connectivity}
Bridge maintenance backlogs and uneven broadband coverage are classic
\textbf{infrastructure} themes. The fixture encodes both physical asset
risk and digital divide geography.

Handbook maintenance teams can prioritize bond packages or grant
applications using gap flags when coverage scores fall below policy
thresholds.

\textbf{Evidence rows.} \texttt{ev-007} (weight 0.9, DOT condition
database export) states deferred bridge maintenance spans 47 structures
flagged as priority repair---high confidence because asset registries
are auditable. \texttt{ev-008} (weight 0.75, Broadband map consortium)
reports fiber-to-premises coverage reaches 78\% of households with rural
gaps along the northern ridge. Together they pair heavy civil liability
with last-mile equity, a common pairing in regional handbooks.

\textbf{Synthesis and finance.} Strong weights on both rows help the
infrastructure section score stay robust under default thresholds. If
bond scoring workflows consume \texttt{handbook\_report.json}, teams can
watch \texttt{scores\_by\_section} for \texttt{07\_infrastructure}
alongside raw evidence counts to see whether narrative urgency matches
quantitative backing.

\textbf{Digital divide nuance.} The broadband row is explicitly
geographic (``northern ridge''), which handbook prose should preserve
rather than generalizing to ``rural areas'' without naming the ridge
feature, unless new evidence broadens the claim.

\textbf{Operations.} When DOT exports a new quarterly extract, version
the corpus and add or update rows so \texttt{reviewed\_at} reflects the
extract date {[}@template2026{]}.
\end{frame}

\end{document}
